% ======================================================================
\section{Preliminaries}\label{sec:prelim}
% ======================================================================

This section collects every tool used later. Each fact comes with a short proof, so that
the main arguments can be checked line by line.

\subsection{Notation}

Vectors live in $\R^d$. For $a,b\in\R^d$, $\ip{a}{b}=\sum_{i=1}^d a_ib_i$ is the inner
product and $\norm{a}=\sqrt{\ip{a}{a}}$ the Euclidean norm. For $f:\R^d\to\R$ we write
$\fs:=\inf_{x\in\R^d}f(x)$ (possibly not attained). Sums $\sum_{k<K}$ run over
$k=0,1,\dots,K-1$. For SGDM we will use the following names, all defined precisely in
\cref{sec:momentum}:
\begin{center}\small
\begin{tabular}{@{}ll@{}}
\toprule
$\alpha>0$, $\beta\in[0,1)$ & step size and momentum parameter\\
$\gamma=\alpha/(1-\beta)$ & effective step size\\
$v_k=x_k-x_{k-1}$ & velocity (last displacement), $v_0=0$\\
$z_k=x_k+\frac{\beta}{1-\beta}v_k$ & extrapolated point\\
$\cb=\frac{L\beta^2}{2(1-\beta)^2}$ & weight of the kinetic energy\\
$\Phi_k=f(z_k)-\fs+\cb\norm{v_k}^2$ & potential\\
$\eps_k=g_k-\nabla f(x_k)$ & gradient noise\\
$\Dlt=f(x_0)-\fs$ & initial optimality gap\\
\bottomrule
\end{tabular}
\end{center}

\subsection{Elementary inequalities}

\begin{lemma}[Elementary facts]\label{lem:elementary}
For all $a,b\in\R^d$:
\begin{enumerate}[label=\textup{(\alph*)}]
  \item \emph{(Cauchy--Schwarz)} $|\ip{a}{b}|\le\norm{a}\norm{b}$.
  \item \emph{(Young)} For every $\rho>0$:
        $\ip{a}{b}\le\frac{\rho}{2}\norm{a}^2+\frac{1}{2\rho}\norm{b}^2$.
  \item \emph{(Polarization)} $-\ip{a}{b}=\frac12\norm{a-b}^2-\frac12\norm{a}^2-\frac12\norm{b}^2$.
  \item \emph{(Convexity of $\norm{\cdot}^2$)} For every $\theta\in[0,1]$:
        $\norm{\theta a+(1-\theta)b}^2\le\theta\norm{a}^2+(1-\theta)\norm{b}^2$.
  \item $\norm{a+b}^2\le2\norm{a}^2+2\norm{b}^2$.
  \item \emph{(Geometric series)} For $q\in[0,1)$ and $k\ge1$:
        $\sum_{j=0}^{k-1}q^j=\frac{1-q^k}{1-q}\le\frac{1}{1-q}$.
\end{enumerate}
\end{lemma}

\begin{proof}
(a) is the Cauchy--Schwarz inequality from linear algebra.
(b) $0\le\Norm{\sqrt\rho\,a-b/\sqrt\rho}^2=\rho\norm{a}^2-2\ip{a}{b}+\norm{b}^2/\rho$;
rearrange and divide by $2$.
(c) Expand $\norm{a-b}^2=\norm{a}^2-2\ip{a}{b}+\norm{b}^2$ and rearrange.
(d) Expanding both sides,
\[
\theta\norm{a}^2+(1-\theta)\norm{b}^2-\norm{\theta a+(1-\theta)b}^2
=\theta(1-\theta)\bigl(\norm{a}^2-2\ip{a}{b}+\norm{b}^2\bigr)=\theta(1-\theta)\norm{a-b}^2\ge0 .
\]
(e) Apply (d) with $\theta=\frac12$ to the vectors $2a$ and $2b$:
$\norm{a+b}^2=\norm{\tfrac12(2a)+\tfrac12(2b)}^2\le\tfrac12\norm{2a}^2+\tfrac12\norm{2b}^2$.
(f) Multiply $\sum_{j=0}^{k-1}q^j$ by $1-q$; the sum telescopes to $1-q^k\le1$.
\end{proof}

\subsection{Smooth functions and the descent lemma}

\begin{definition}[$L$-smoothness]\label{def:smooth}
A differentiable $f:\R^d\to\R$ is \emph{$L$-smooth} ($L>0$) if its gradient is
$L$-Lipschitz:
$\norm{\nabla f(x)-\nabla f(y)}\le L\norm{x-y}$ for all $x,y\in\R^d$.
\end{definition}

For twice differentiable $f$, $L$-smoothness is equivalent to all eigenvalues of the
Hessian $\nabla^2 f(x)$ lying in $[-L,L]$; for example $f(x)=\frac h2 x^2$ is $|h|$-smooth.
The single most important consequence is that an $L$-smooth function lies below a
parabola around every point.

\begin{lemma}[Descent lemma]\label{lem:descent}
If $f$ is $L$-smooth, then for all $x,y\in\R^d$,
\begin{equation}\label{eq:descent}
  f(y)\le f(x)+\ip{\nabla f(x)}{y-x}+\frac L2\norm{y-x}^2 .
\end{equation}
\end{lemma}

\begin{proof}
Let $\varphi(t)=f(x+t(y-x))$ for $t\in[0,1]$. By the chain rule
$\varphi'(t)=\ip{\nabla f(x+t(y-x))}{y-x}$, which is continuous because $\nabla f$ is
Lipschitz. By the fundamental theorem of calculus,
$f(y)-f(x)=\varphi(1)-\varphi(0)=\int_0^1\varphi'(t)\,dt$. Adding and subtracting
$\ip{\nabla f(x)}{y-x}$ inside the integral,
\begin{align*}
f(y)-f(x)
&=\ip{\nabla f(x)}{y-x}+\int_0^1\ip{\nabla f(x+t(y-x))-\nabla f(x)}{y-x}\,dt\\
&\le\ip{\nabla f(x)}{y-x}+\int_0^1\norm{\nabla f(x+t(y-x))-\nabla f(x)}\,\norm{y-x}\,dt\\
&\le\ip{\nabla f(x)}{y-x}+\int_0^1 L\,t\norm{y-x}^2\,dt
 =\ip{\nabla f(x)}{y-x}+\frac L2\norm{y-x}^2,
\end{align*}
where we used Cauchy--Schwarz (\cref{lem:elementary}(a)) and then $L$-smoothness.
\end{proof}

\begin{lemma}[Gradients are controlled by the optimality gap]\label{lem:grad-gap}
If $f$ is $L$-smooth and $\fs>-\infty$, then
$\norm{\nabla f(x)}^2\le 2L\bigl(f(x)-\fs\bigr)$ for every $x\in\R^d$.
\end{lemma}

\begin{proof}
Apply \cref{lem:descent} with $y=x-\frac1L\nabla f(x)$:
$\fs\le f(y)\le f(x)-\frac1L\norm{\nabla f(x)}^2+\frac1{2L}\norm{\nabla f(x)}^2
=f(x)-\frac1{2L}\norm{\nabla f(x)}^2$. Rearrange.
\end{proof}

\subsection{Convexity and the Polyak--{\L}ojasiewicz inequality}

\begin{definition}\label{def:convex}
A differentiable $f$ is \emph{convex} if $f(y)\ge f(x)+\ip{\nabla f(x)}{y-x}$ for all
$x,y$, and \emph{$\mu$-strongly convex} ($\mu>0$) if
$f(y)\ge f(x)+\ip{\nabla f(x)}{y-x}+\frac\mu2\norm{y-x}^2$ for all $x,y$.
\end{definition}

(For differentiable functions these first-order inequalities are equivalent to the usual
definitions of convexity; we take them as definitions.)

\begin{definition}[PL inequality]\label{def:pl}
An $f$ with $\fs>-\infty$ satisfies the \emph{Polyak--{\L}ojasiewicz (PL) inequality} with
constant $\mu>0$ if
\begin{equation}\label{eq:pl}
  \norm{\nabla f(x)}^2\ge2\mu\bigl(f(x)-\fs\bigr)\qquad\text{for all }x\in\R^d.
\end{equation}
\end{definition}

The PL inequality says that the gradient can be small only where the function value is
nearly optimal; in particular every stationary point is a global minimizer. It does
\emph{not} require convexity.

\begin{lemma}\label{lem:pl-facts}
\begin{enumerate}[label=\textup{(\alph*)}]
\item A $\mu$-strongly convex function satisfies the PL inequality with constant $\mu$.
\item If $f$ is $L$-smooth and satisfies PL with constant $\mu$, and $\nabla f$ is not
      identically zero, then $\mu\le L$.
\item The function $f(x)=x^2+3\sin^2x$ on $\R$ is nonconvex, $8$-smooth, and satisfies PL
      (with $\mu=1/32$, see~\cite{karimi2016}); so does $\sum_{i=1}^d(x_i^2+3\sin^2x_i)$ on $\R^d$.
\end{enumerate}
\end{lemma}

\begin{proof}
(a) Fix $x$. The right-hand side of the strong-convexity inequality is a quadratic
function of $y$, minimized at $y=x-\frac1\mu\nabla f(x)$ with minimum value
$f(x)-\frac{1}{2\mu}\norm{\nabla f(x)}^2$. Hence $f(y)\ge f(x)-\frac1{2\mu}\norm{\nabla
f(x)}^2$ for every $y$; taking the infimum over $y$ gives
$\fs\ge f(x)-\frac1{2\mu}\norm{\nabla f(x)}^2$, which is \eqref{eq:pl}.

(b) Pick $x$ with $\nabla f(x)\ne0$. By \cref{lem:grad-gap} and \eqref{eq:pl},
$\frac{1}{2L}\norm{\nabla f(x)}^2\le f(x)-\fs\le\frac1{2\mu}\norm{\nabla f(x)}^2$; divide by
$\norm{\nabla f(x)}^2>0$.

(c) $f''(x)=2+6\cos2x\in[-4,8]$, so $f$ is $8$-smooth and (as $f''$ takes negative values)
nonconvex. The PL constant is computed in~\cite{karimi2016}; for a sum of functions of
separate coordinates, \eqref{eq:pl} holds coordinate-wise and therefore for the sum.
\end{proof}

\subsection{Randomness and conditional expectation}\label{sec:conditional}

\paragraph{The probabilistic model.} The samples $\xi_0,\xi_1,\xi_2,\dots$ are independent
random variables, and the starting point $x_0$ is deterministic. The stochastic gradient
oracle $g(x;\xi)$ satisfies, for every \emph{fixed} $x\in\R^d$,
\begin{equation}\label{eq:oracle}
  \E_\xi\bigl[g(x;\xi)\bigr]=\nabla f(x),\qquad
  \E_\xi\Norm{g(x;\xi)-\nabla f(x)}^2\le\sigma^2 .
\end{equation}

\paragraph{What is known at time $k$.} Every quantity computed by the algorithm
\emph{before} it draws $\xi_k$---the points $x_1,\dots,x_k$, and hence $\nabla f(x_k)$,
$f(x_k)$, $v_k=x_k-x_{k-1}$, $z_k$, $\Phi_k$---is a function of $\xi_0,\dots,\xi_{k-1}$.
We say such a quantity is \emph{known at time $k$}. The quantities $g_k$, $x_{k+1}$,
$z_{k+1}$, $\Phi_{k+1}$ are not known at time $k$ (they depend on $\xi_k$). We write
\[
  \Ek[Y]:=\E\bigl[Y\,\big|\,\xi_0,\dots,\xi_{k-1}\bigr]
\]
for the conditional expectation of a random variable $Y$ given the history up to time $k$.
Intuitively, $\Ek[Y]$ averages $Y$ over the samples that have not been drawn yet, holding
the history fixed; it is itself a random quantity (a function of the history). We use only
the following standard rules (valid whenever the expectations involved are finite):
\begin{enumerate}[label=\textup{(R\arabic*)},leftmargin=3.2em]
\item\label{R:lin} \emph{Linearity:} $\Ek[aY+bZ]=a\Ek[Y]+b\Ek[Z]$ for constants $a,b$.
\item\label{R:known} \emph{Known factors come out:} if $U$ is known at time $k$ then
  $\Ek[U\,Y]=U\,\Ek[Y]$ and $\Ek[U]=U$; for vectors, $\Ek\ip{U}{Y}=\ip{U}{\Ek[Y]}$.
\item\label{R:tower} \emph{Tower property:} $\E\bigl[\Ek[Y]\bigr]=\E[Y]$.
\item\label{R:mono} \emph{Monotonicity:} if $Y\le Z$ then $\Ek[Y]\le\Ek[Z]$.
\item\label{R:freeze} \emph{Freezing:} if $U$ is known at time $k$, then
  $\Ek\bigl[h(U,\xi_k)\bigr]=H(U)$, where $H(u):=\E_\xi[h(u,\xi)]$ is the ordinary
  expectation computed with $u$ held fixed. (This uses that $\xi_k$ is independent of
  $\xi_0,\dots,\xi_{k-1}$.)
\end{enumerate}
Applying \ref{R:freeze} to \eqref{eq:oracle} with $U=x_k$ gives the two properties that
the analysis actually uses:
\begin{equation}\label{eq:cond-oracle}
  \Ek[g_k]=\nabla f(x_k),\qquad \Ek\norm{\eps_k}^2\le\sigma^2,
  \qquad\text{where }\eps_k:=g_k-\nabla f(x_k).
\end{equation}

\begin{lemma}[Variance decomposition]\label{lem:variance}
Let $Y$ be a random vector with $\Ek\norm{Y}^2<\infty$ whose conditional mean $m:=\Ek[Y]$
is known at time $k$ (as it always is). Then $\Ek\norm{Y}^2=\norm{m}^2+\Ek\norm{Y-m}^2$.
In particular $\Ek\norm{g_k}^2=\norm{\nabla f(x_k)}^2+\Ek\norm{\eps_k}^2
\le\norm{\nabla f(x_k)}^2+\sigma^2$.
\end{lemma}

\begin{proof}
$\norm{Y}^2=\norm{m+(Y-m)}^2=\norm{m}^2+2\ip{m}{Y-m}+\norm{Y-m}^2$. Take $\Ek$: by
\ref{R:known}, $\Ek\ip{m}{Y-m}=\ip{m}{\Ek[Y]-m}=0$. The second claim is the case $Y=g_k$,
$m=\nabla f(x_k)$, combined with \eqref{eq:cond-oracle}.
\end{proof}

\begin{remark}[Mini-batches]\label{rem:minibatch}
If $g_k$ is the average of $B$ stochastic gradients computed with independent samples at
the same point $x_k$, then \eqref{eq:cond-oracle} holds with $\sigma^2$ replaced by
$\sigma^2/B$: the noises are independent with mean zero, so all cross terms in
$\Ek\norm{\frac1B\sum_{b=1}^B\eps_{k,b}}^2$ vanish and $B$ terms of size at most
$\sigma^2/B^2$ remain.
\end{remark}

\subsection{Standing assumptions}

\begin{assumption}[Smoothness]\label{as:smooth}
$f$ is $L$-smooth (\cref{def:smooth}).
\end{assumption}
\begin{assumption}[Lower bound]\label{as:lower}
$\fs=\inf_x f(x)>-\infty$.
\end{assumption}
\begin{assumption}[Unbiased gradients]\label{as:unbiased}
$\Ek[g_k]=\nabla f(x_k)$ for every $k\ge0$.
\end{assumption}
\begin{assumption}[Bounded variance]\label{as:variance}
$\Ek\norm{g_k-\nabla f(x_k)}^2\le\sigma^2$ for every $k\ge0$.
\end{assumption}

\Cref{as:unbiased,as:variance} hold under the model \eqref{eq:oracle} by \eqref{eq:cond-oracle};
the proofs below use only these conditional versions, so they also cover other sampling
schemes for which \cref{as:unbiased,as:variance} hold. The case $\sigma=0$ is
deterministic (exact gradients).

\begin{remark}[All expectations are finite]\label{rem:integrability}
Under \cref{as:smooth,as:variance}, every quantity whose expectation we take is integrable.
Indeed, by induction on $k$ one has $\E\norm{x_k}^2<\infty$: if this holds up to $k$,
then $\E\norm{g_k}^2=\E\bigl[\Ek\norm{g_k}^2\bigr]\le\E\norm{\nabla f(x_k)}^2+\sigma^2
\le2\norm{\nabla f(0)}^2+2L^2\E\norm{x_k}^2+\sigma^2<\infty$ (\cref{lem:variance},
\ref{R:tower}, $L$-smoothness and \cref{lem:elementary}(e)), and $x_{k+1}$ is a linear
combination of $x_k$, $x_{k-1}$ and $g_k$. Finally
$\fs\le f(x)\le f(0)+\norm{\nabla f(0)}\norm{x}+\frac L2\norm{x}^2$ by \cref{lem:descent}
and Cauchy--Schwarz, so $\E|f(x_k)|<\infty$, and similarly for $z_k$ and $\Phi_k$. We
will not mention integrability again.
\end{remark}
